In contrast to simple \LaTeX documents, which usually consist of only one
source file, this template consists of several files and subdirectories. This
makes it easier to structure the content and separate it from the layout
adjustments. The directory structure may seem a little opaque at first, but it
will become clearer after the following explanations.
Figure~\ref{fig:tpl:dir} represents the directory structure
tree-like.

\begin{figure}
\dirtree{%
.1 /.
.2 thesis.tex\DTcomment{main file}.
.2 latex/\DTcomment{subdirectory for \LaTeX customisations}.
.3 float.tex\DTcomment{Standard settings for floating environments}.
.3 titlepage.tex\DTcomment{Format templates for cover page}.
.3 \dots.
.2 figures/\DTcomment{subdirectory for illustrations}.
.3 tuc\_black.pdf\DTcomment{TUC logo black}.
.3 tuc\_black\_margin.pdf\DTcomment{TUC logo black with margin}.
.3 tuc\_green.pdf\DTcomment{TUC logo green}.
.3 tuc\_green\_margin.pdf\DTcomment{TUC logo green with border}.
.3 \dots.
.2 content/\DTcomment{content subdirectory}.
.3 titlepage.tex\DTcomment{Cover page}.
.3 \dots.
.2 bibliogrpahy.bib\DTcomment{literature database}.
.2 Makefile\DTcomment{control file for automatic \LaTeX run using \texttt{make}}.
}
\caption{Template directory structure}.
\label{fig:tpl:dir}
\end{figure}

\lstinputlisting[language={[LaTeX]{TeX}},style=block,float={p},caption={Principal structure of the main file \texttt{thesis.tex}},label={lst:tpl:thesis},morekeywords={chapter,subsection}]{code/tpl_thesis.en.tex}

The file \texttt{thesis.tex} is the actual main file created by the
\LaTeX compiler. In it the document class is defined and additional packages
are loaded. Listing~\ref{lst:tpl:thesis} shows the basic structure of this
file. In the template, this file is more extensive and supplemented by comments.
After loading the additional packages, further files from the subdirectory
\texttt{latex} are included. These make adjustments to the standard settings or define
further commands. Within the file \texttt{float.tex} e.g. the placement settings for
floating environments are made. Within this directory you can also store your
own adjustments. This keeps the main file clear.

In the directory~\texttt{figures} you should store all illustrations for your work.
If necessary, you can structure the directory yourself by adding further
subdirectories. The TUC logos are already in this directory, as they are
needed, for example, for the design of the cover page.

No output has been created yet. In the directory~\texttt{content} you can store the
content of your paper. Structure at least on chapter level -- better on section
level. Within the main file only the outline commands should appear. You can
then use \verb+\input{}+ after the outline command to include the respective
content.

For authors who prefer to work on the command line, the file \texttt{Makefile} is
helpful. It enables the simple translation of the project with the utility
\texttt{make}. If you have never worked with this programme, simply ignore this file.
It has no no relevance for you.

The file~\texttt{literature.bib} in the main directory contains the literature database.
In its initial state, the template does not generate a bibliography.
It must first be set up according to section~\ref{sec:content:bibliography} so
that bibliographies are automatically generated from this literature database.
